% =============================================================================
% FORMALIA-SCHALTER -- hier stellst du deinen Leitfaden ein
% =============================================================================
% An der FOM gibt es nicht DEN einen Leitfaden. Welcher fuer dich gilt, steht
% in deiner Modulbeschreibung bzw. sagen dir deine Pruefenden. Stelle hier ein,
% welcher gilt -- die Vorlage setzt Raender, Schrift, Zitierstil und die
% leitfadenspezifischen Sonderregeln dann automatisch richtig.
%
% Hintergrund zu jeder Einstellung: ../anleitung/02_formalia.md
%
% WICHTIG: Individuelle Vorgaben deiner Pruefenden haben IMMER Vorrang vor dem
% Leitfaden. Solche Abweichungen traegst du nicht hier ein, sondern in
% formalia/profil_eigen.tex -- dann bleibt nachvollziehbar, was Leitfaden ist
% und was Pruefervorgabe.
% =============================================================================

% --- Welcher Leitfaden gilt? -------------------------------------------------
%   wi   = Gestaltung wissenschaftlicher Arbeiten, V1.4 (03/2024)
%          Wirtschaftsinformatik & Ingenieurwesen
%   jks  = Formale Gestaltung, Jaeger/Kuempel/Seng (01/2024)
%          allgemeiner FOM-Standard (Wirtschaft & Recht u. a.)
\def\FormaliaLeitfaden{jks}

% --- Zitierstil --------------------------------------------------------------
%   fussnote    = Kurzbeleg in der Fussnote  (JKS-Standard, WI zulaessig)
%   authoryear  = Harvard/Autor-Jahr im Text (JKS zulaessig, WI zulaessig)
%   ieee        = numerisch [12]             (nur WI, in der Informatik ueblich)
%
% ACHTUNG: Diesen Wert einmal am Anfang festlegen und dann NICHT mehr wechseln.
% Der Schalter aendert nur die Ausgabe, nicht den Sinn der bereits im Text
% gesetzten \autocite-Befehle. Ein spaeter Wechsel von fussnote auf authoryear
% verschiebt saemtliche Belege aus den Fussnoten in den Fliesstext -- das musst
% du dann sprachlich nacharbeiten.
\def\FormaliaZitierstil{fussnote}

% --- Schrift -----------------------------------------------------------------
%   tnr    = Times New Roman 12 pt (Standard beider Leitfaeden)
%            Gesetzt wird TeX Gyre Termes, ein metrisch identischer freier Klon.
%   arial  = Arial als zulaessige Alternative (WI 11 pt / JKS 11,5 pt)
%            Gesetzt wird TeX Gyre Heros, ein metrisch identischer freier Klon.
\def\FormaliaSchrift{tnr}

% --- Optionales Overlay ------------------------------------------------------
%   ifes   = zusaetzliche Vorgaben fuer empirische Arbeiten (Gansser u. a.)
%   kein   = kein Overlay
% Das Overlay ersetzt den Leitfaden nicht, es ergaenzt ihn.
\def\FormaliaOverlay{kein}
